% Proposed strengthening-stage text. NOT automatically included in main.tex.
% No new empirical claim is made by this file.
\section{Diagnosing the Opportunity for Selective Intervention}

An authority-limited controller should not be assessed solely by the number of
agents or treatments. The evaluation must establish whether the permitted controls
can materially alter harm, and whether evidence-informed selection can improve upon
uninformed allocation under those same limits. We separate fixed-proposal
accounting, diagnostic control opportunity, and prospective closed-loop evaluation.

\paragraph{Fixed-proposal accounting.}
For a recorded sequence of per-agent attempted violation indicators
$v_{i,t}\in\{0,1\}$, assume a deterministic gateway can eliminate that declared
per-agent violation when agent $i$ is selected, without changing other agents'
recorded proposals. For disjoint, prespecified intervention windows $W_j$, let
$b_{j,i}=\sum_{t\in W_j}v_{i,t}$. If precisely $k$ of $N$ agents are selected in
each window, the largest direct reduction on this fixed record is
\[
 G^*_{\mathrm{record}}=\frac{1}{NT}\sum_j
          \sum_{i\in\operatorname{Top}_k(b_j)} b_{j,i}.
\]
Uniform random selection has expected direct reduction
\[
 \mathbb E[G_{\mathrm{random}}\mid(v_{i,t})]
 =\frac{1}{NT}\sum_j\frac{k}{N}\sum_i b_{j,i}.
\]
Their difference quantifies the available selection opportunity for this record
and schedule. When all agents have identical window violation counts, this
difference is zero: even hindsight selection cannot improve on uniform selection
on the recorded per-agent endpoint. Unequal counts create potential opportunity,
not a guarantee that a causal online ranker can exploit it.

\paragraph{Proof and scope.}
Each window's reduction is additive over selected agents, so selecting its $k$
largest counts maximizes that window's reduction. Disjoint windows permit summing
these maxima. Each identity has inclusion probability $k/N$ under a uniform subset,
which gives the random-selection expectation by linearity, without independence
between agents. The result is a hindsight diagnostic, not an implementable
forecaster, not an optimum over intervention schedules, and not a bound on changes
to future learning or physical joint events. All direct reductions are bounded by
realized authorized agent-time on the same record. The earlier long-run authority
bound therefore remains relevant, but cannot establish that a particular controller
allocates its authority optimally.

\paragraph{Independent service utility.}
For the original saturated resource-sharing setting, averaging
$r_i=x_i-0.2v_i+0.3\bar x$ over $N=50$ agents with total allocation $100$ gives
$\bar r=2.6-0.2\bar v$. Its reward improvement is not an independent measure of
legitimate demand satisfaction or throughput. Those service claims require
separately specified task measures; they are not inferred from this identity.

\paragraph{Timing and predictive-control diagnostics.}
The new development instrumentation records the complete sequence from observed
risk through alarm, admission, command modification, and contact. An alarm at the
end of step $t$ normally authorizes an intervention at $t+1$. The separately named
zero-delay state-model diagnostic acts before the current transition and changes
several information/timing assumptions together; it is not a clean single-factor
ablation. Missing events remain absent/censored rather than being assigned zero
latency. Whole-episode contact incidence is retained alongside prospectively
specified timing and duration summaries.

A candidate predictive primitive evaluates bounded commands using an eight-step
constant-command, contact-free dynamics approximation with the same actuator lag.
All compared triggers and selectors share that primitive within its block. A
state-model benefit score estimates single-target reduction of that model's risk;
it neither observes future simulator states nor establishes causal culpability.
A reduction of model cost is not empirical evidence of certified physical safety.
One-step free-motion agreement and realized prediction errors are exported to expose
where the approximation holds and fails.

\paragraph{Development versus confirmation.}
These diagnostics retain all development seeds and conditions. They do not supply
confirmatory significance or select a favorable result for the manuscript. A future
confirmation must freeze the chosen method, comparable calibration opportunities,
authority constraints, task requirements and primary contrasts before new independent
policy seeds are evaluated. Existing null and unfavorable results remain preserved.
